\documentclass[11pt,a4paper,oneside]{book}
\usepackage[utf8]{inputenc}
\usepackage[T1]{fontenc}
\usepackage{lmodern}
\usepackage[margin=2.3cm,headheight=14pt]{geometry}
\usepackage{amsmath,amssymb,amsfonts,mathtools,bm}
\usepackage{longtable,booktabs,array,calc}
\usepackage{float}
\usepackage[dvipsnames]{xcolor}
\usepackage{tikz}
\usetikzlibrary{trees,positioning,arrows.meta}
\usepackage{pgfplots}
\pgfplotsset{compat=1.16}
\usepgfplotslibrary{fillbetween}
\usepackage[most]{tcolorbox}
\usepackage{enumitem}
\usepackage{fancyhdr}
\usepackage{titlesec}
\usepackage{caption}
\usepackage[hidelinks,bookmarksnumbered]{hyperref}
\usepackage{pifont}
\usepackage{multicol}
\usepackage{tocloft}
\renewcommand{\cftchapfont}{\bfseries\small}\renewcommand{\cftchappagefont}{\bfseries\small}
\renewcommand{\cftsecfont}{\scriptsize}\renewcommand{\cftsecpagefont}{\scriptsize}
\setlength{\cftbeforechapskip}{3pt}\setlength{\cftbeforesecskip}{0pt}
\setlength{\cftsecindent}{0.8em}\setlength{\cftsecnumwidth}{0pt}\setlength{\cftchapnumwidth}{0pt}
\renewcommand{\cftsecleader}{\cftdotfill{6}}\renewcommand{\cftchapleader}{\cftdotfill{6}}
\renewcommand{\cfttoctitlefont}{\Large\bfseries\color{chapcol}}
\setlength{\cftaftertoctitleskip}{6pt}
\renewcommand{\checkmark}{\ding{51}}

% pandoc helpers
\providecommand{\tightlist}{\setlength{\itemsep}{0pt}\setlength{\parskip}{0pt}}
\newcommand{\passthrough}[1]{#1}
\setlength{\parindent}{0pt}
\setlength{\parskip}{5pt plus 2pt}
\setlist{itemsep=2pt,topsep=3pt}
\allowdisplaybreaks

% boxes
\newtcolorbox{exambox}[1]{enhanced,breakable,colback=red!3!white,colframe=red!65!black,
  fonttitle=\bfseries\small,title={$\bigstar$ EXAM --- #1},left=6pt,right=6pt,top=4pt,bottom=4pt,before skip=8pt,after skip=8pt}
\newtcolorbox{infobox}{enhanced,breakable,colback=blue!3!white,colframe=blue!50!black,
  left=6pt,right=6pt,top=4pt,bottom=4pt,before skip=8pt,after skip=8pt}

% headings
\definecolor{chapcol}{RGB}{20,50,110}
\titleformat{\chapter}[display]{\normalfont\huge\bfseries\color{chapcol}}{}{0pt}{\Huge}[\vspace{2pt}{\color{chapcol}\titlerule[1.2pt]}]
\titlespacing*{\chapter}{0pt}{-20pt}{18pt}
\titleformat{\section}{\normalfont\Large\bfseries\color{chapcol}}{}{0pt}{}
\titleformat{\subsection}{\normalfont\large\bfseries\color{chapcol!80!black}}{}{0pt}{}
\titleformat{\subsubsection}{\normalfont\normalsize\bfseries}{}{0pt}{}
\setcounter{secnumdepth}{-1}
\renewcommand{\thefigure}{\arabic{figure}}
\setcounter{tocdepth}{1}

\pagestyle{fancy}
\fancyhf{}
\fancyhead[L]{\small\leftmark}
\fancyhead[R]{}
\fancyfoot[C]{\thepage}
\renewcommand{\chaptermark}[1]{\markboth{#1}{}}
\fancyfoot[R]{\scriptsize 20188 Quant Finance Bible}

\begin{document}
\begin{titlepage}
\centering
\vspace*{3cm}
{\color{chapcol}\rule{\textwidth}{1.5pt}}\\[0.6cm]
{\Huge\bfseries\color{chapcol} The Quant Finance Bible\\[0.4cm]}
{\Large 20188 Quantitative Finance and Derivatives --- Module 1\\[0.2cm]
Arbitrage Theory in Discrete and Continuous Time}\\[0.4cm]
{\color{chapcol}\rule{\textwidth}{1.5pt}}\\[1.5cm]
{\large Theory from the lecture slides (with the material given orally added),\\
the complete official list of exam theory topics 2026--27,\\
every past exam question 2022--2026 with step-by-step solutions,\\
and a final formula sheet.}\\[2cm]
\begin{tikzpicture}
\begin{axis}[width=9cm,height=4.5cm,axis lines=none,xmin=0,xmax=1]
\addplot[chapcol,thick] coordinates {(0,1)(0.05,1.03)(0.1,0.99)(0.15,1.06)(0.2,1.04)(0.25,1.11)(0.3,1.08)(0.35,1.15)(0.4,1.12)(0.45,1.2)(0.5,1.17)(0.55,1.25)(0.6,1.21)(0.65,1.3)(0.7,1.27)(0.75,1.36)(0.8,1.33)(0.85,1.42)(0.9,1.38)(0.95,1.47)(1,1.45)};
\end{axis}
\end{tikzpicture}
\vfill
{\small Reference textbook: [BOR26] Battauz, A., F. Ortu and F. Rotondi (2026), \emph{Arbitrage Theory in Discrete and Continuous Time}.}
\end{titlepage}
\begin{multicols}{2}
\tableofcontents
\end{multicols}
\clearpage
\markboth{How to Use This Bible}{}
\input{build/ch00.tex}
\input{build/ch01.tex}
\input{build/ch02.tex}
\input{build/ch03_04.tex}
\input{build/ch05.tex}
\input{build/ch07.tex}
\input{build/ch08.tex}
\input{build/ch09.tex}
\input{build/ch10.tex}
\input{build/ch11.tex}
\input{build/ch14.tex}
\input{build/ch15.tex}
\input{build/ch16.tex}
\input{build/ch99.tex}
\end{document}
